\chapter{The Cosmic Bit Budget: Spacetime Resolution and the 184-Bit Universe}

In the previous chapter we watched a universe assemble itself from nothing
but an entropy-increasing bitstring. Expansion, matter tiers, and the arrow
of time all emerged without a single hardcoded physical law. But qualitative
emergence is not enough. A serious theory must produce numbers --- specific,
checkable predictions that can be compared against observation.

This chapter asks a precise question: if the universe operates on a finite
bit budget of $n$ bits, what is the exact resolution of space and time? The
answer will be a single integer. And that integer, derived from one
measurement, will correctly predict two completely independent results from
cosmology and black hole physics.

\section{Two Resolutions from One Number}

A universe of $n$ bits has exactly $2^n$ distinct configurations. If space
is the geometric reading of informational variety, then $2^n$ is the maximum
number of distinguishable spatial positions the system can represent. We
identify each position with one Planck length --- the smallest length scale
physics can define --- giving a maximum radial extent of

\[
\Delta x_{\max} = 2^n \cdot \ell_P.
\]

Time requires a different argument. Time, as established in Chapter~1, is
not an external parameter but an internal ordering relation: the sequence of
states an observer traverses as the system evolves from its zero-entropy
initial state toward thermodynamic equilibrium. The natural unit of time is
therefore not an arbitrary clock tick but the number of bit-flip steps
required to bring a zero-entropy system to full saturation.

This is a classical result from combinatorics known as the coupon-collector
problem. If you flip bits one at a time, randomly and uniformly, it takes on
average

\[
n \ln n
\]

steps to visit every one of the $n$ bit positions at least once --- that is,
to reach full statistical equilibrium. Identifying one bit-flip with one
Planck time:

\[
\Delta t_{\max} = n \ln n \cdot t_P.
\]

These two resolutions --- one exponential in $n$, one quasi-linear --- define
the geometry of the spacetime block. Their ratio is the key quantity.

\section{The Aspect Ratio}

Dividing the spatial resolution by the temporal resolution gives a
dimensionless number:

\[
\mathcal{A}(n) = \frac{2^n}{n \ln n}.
\]

We call this the \emph{aspect ratio} of the spacetime block. It is not a
free parameter. Given $n$, it is fixed. It encodes the complete geometric
relationship between the spatial and temporal extent of the universe from a
single integer.

Because $2^n$ grows exponentially while $n \ln n$ grows quasi-linearly,
the aspect ratio is strictly increasing: larger bit budgets produce
universes where space is exponentially larger relative to time. This
monotonicity has an immediate corollary. Any subsystem of the universe ---
a black hole, a galaxy, a particle --- has a smaller bit budget than the
whole, and therefore a strictly smaller aspect ratio. The hierarchy is
structural, not assumed.

\section{Reading the Inflationary Epoch}

To find $n$ for our universe, we need a physical measurement that can be
identified with $\mathcal{A}(n)$. The inflationary epoch of the early
universe is the natural candidate. Inflation is, physically, the closest
realisation of what the model describes: a system expanding rapidly from a
near-zero-entropy initial state.

Standard cosmology gives the following values for the inflationary epoch:
a duration of approximately $10^{-35}$ seconds and a radial expansion of
approximately $10^{26}$ metres. To compare these against $\mathcal{A}(n)$,
we must strip away human units and express both in Planck units:

\[
\Delta t_P = \frac{10^{-35}}{5.39 \times 10^{-44}} \approx 1.86 \times 10^8
\text{ Planck times,}
\]
\[
\Delta x_P = \frac{10^{26}}{1.616 \times 10^{-35}} \approx 6.19 \times 10^{60}
\text{ Planck lengths.}
\]

The dimensionless inflationary aspect ratio is then:

\[
\mathcal{A}_{\mathrm{inf}} = \frac{6.19 \times 10^{60}}{1.86 \times 10^8}
\approx 3.3 \times 10^{52}.
\]

Setting $\mathcal{A}(n) = \mathcal{A}_{\mathrm{inf}}$ and solving numerically:

\[
\frac{2^n}{n \ln n} = 3.3 \times 10^{52}
\quad \Longrightarrow \quad
\boxed{n \approx 184 \text{ bits.}}
\]

One measurement. One integer. That is the universe's bit budget.

\section{An Independent Consistency Check}

Standard inflationary cosmology requires, independently of anything above,
that the early universe underwent a minimum of 60 e-folds of exponential
expansion --- about 87 doublings on a base-2 scale --- in order to explain
why the universe looks so flat and uniform today. This lower bound comes
from a completely different line of reasoning: the horizon and flatness
problems of classical cosmology.

A bit budget of $n \approx 184$ corresponds to approximately 127 natural-log
units of expansion capacity.

We did not use the e-fold count to fix $n$. We fixed $n$ from the Planck-unit
ratio of inflationary spatial and temporal scales. The fact that the result
comfortably exceeds the independent cosmological lower bound of 60 e-folds ---
with 127 natural-log units, or about 183 base-2 doublings of spatial depth
--- is a non-trivial consistency check from a completely separate branch of
physics.

\section{Black Holes: A Local Bit Budget}

If $n \approx 184$ describes the full universe, what happens when we apply
the same framework to a smaller, localised system? A black hole is the
natural test case: it is a zero-entropy sink (the singularity, from
Chapter~2) surrounded by an information boundary (the event horizon). Its
characteristic length scale is the Schwarzschild radius $r_s$.

By direct analogy with the universe, whose spatial resolution $2^n$ was
identified with a radial length in Planck units, we assign the black hole
a local bit budget:

\[
n_{\mathrm{bh}} = \log_2\!\left(\frac{r_s}{\ell_P}\right).
\]

For a solar-mass black hole, $r_s \approx 2950$ metres, giving:

\[
n_{\mathrm{bh}} = \log_2\!\left(\frac{2950}{1.616 \times 10^{-35}}\right)
\approx \log_2(1.83 \times 10^{38}) \approx 127 \text{ bits.}
\]

Note that $127 < 184$, as required by the monotonicity result: a black hole
is a subsystem of the universe and must carry a strictly smaller bit budget.

The event horizon is a two-dimensional surface. By the same logic that gives
$2^n$ as the one-dimensional spatial resolution, the two-dimensional surface
resolution is $2^{2n_{\mathrm{bh}}}$ Planck areas. The predicted physical
surface area is:

\[
A_{\mathrm{model}} = 2^{2n_{\mathrm{bh}}} \cdot \ell_P^2
= (1.83 \times 10^{38})^2 \times (1.616 \times 10^{-35})^2
\approx 8.7 \times 10^6 \text{ m}^2.
\]

The Bekenstein--Hawking formula, derived from semiclassical general
relativity, gives:

\[
A_{\mathrm{BH}} = 4\pi r_s^2 \approx 1.09 \times 10^8 \text{ m}^2.
\]

The ratio between the two predictions is:

\[
\frac{A_{\mathrm{model}}}{A_{\mathrm{BH}}} = \frac{1}{4\pi}.
\]

The discrepancy is exactly $4\pi$ --- the ratio of a flat square of side $r$
to the surface area of a sphere of radius $r$. It is not a failure of the
model. It is the signature of spherical geometry. The event horizon is a
sphere because General Relativity demands it. The information-theoretic
calculation computes a flat raster area; multiplying by $4\pi$ to convert
to a spherical surface recovers Bekenstein--Hawking exactly:

\[
A_{\mathrm{model}} \times 4\pi \approx 1.09 \times 10^8 \text{ m}^2
= A_{\mathrm{BH}}. \quad \checkmark
\]

This $4\pi$ factor is the same geometric conversion that appeared in
Chapter~2 when we discussed holographic consistency. Its reappearance here
is not a coincidence. It is the unique factor demanded whenever a
one-dimensional radial information count is compared to a
three-dimensional spherical observable --- the precise boundary between
the informational and geometric descriptions of the same object.

\section{What $G$, $c$, and $\hbar$ Are}

The derivation above introduced no fundamental constants of physics. Newton's
gravitational constant $G$, the speed of light $c$, and Planck's constant
$\hbar$ appear nowhere in the information-theoretic calculation. They enter
only at the very end, when we translate back into human units --- metres,
seconds, kilograms.

This is not an accident. In the framework, $G$, $c$, and $\hbar$ are unit
conversion factors between the informational description and the human
measurement system. They are not inputs to the physics. They are outputs of
the choice of units. The physics is the bit count $n$. The constants are the
dictionary between $n$ and the numbers on our measuring instruments.

\section{What This Chapter Establishes}

Three results follow from the aspect ratio argument.

First, a single integer $n \approx 184$ encodes the full geometric structure
of spacetime. Spatial resolution, temporal resolution, and their ratio are
all determined by $n$ alone through $\mathcal{A}(n) = 2^n / (n \ln n)$.

Second, the e-fold count of inflation --- derived entirely independently from
the horizon and flatness problems --- is consistent with $n \approx 184$.
The two calculations were done in different branches of physics and agree.

Third, the Bekenstein--Hawking entropy of a solar-mass black hole is
recovered from the local bit budget $n_{\mathrm{bh}} = \log_2(r_s/\ell_P)$,
up to the factor $4\pi$ that converts flat raster geometry to the spherical
geometry demanded by GR. No constants of physics are introduced; they emerge
from the unit conversion at the end.

The next chapter asks what happens when an observer is placed inside this
184-bit spacetime and attempts to describe what they see. The answer is
quantum mechanics.
